% ====================================================================
% ФАЙЛ: tasks/02_mkt/mkt_humidity_bank_part1.tex
% ТЕМА: ВЛАЖНОСТЬ ВОЗДУХА. НАСЫЩЕННЫЕ ПАРЫ
% ПАЧКА 1: ВАРИАНТЫ 1–6
% ====================================================================

% === ЗАДАЧА 1 ===
% Тег: #МКТ #ВЛАЖНОСТЬ #КИМ_24 #ВАРИАНТ_01
\begin{taskbasic}[code={\label{task:mkt_hum_v1_n24}}]
\textbf{Задача 1.} В герметичном сосуде объёмом $V = 80$~л находится водяной пар при температуре $t_1 = 150^{\circ}\text{C}$ и давлении $p_1 = 8$~кПа. Какая масса воды $\Delta m$ сконденсируется в сосуде при охлаждении пара до температуры $t_2 = 20^{\circ}\text{C}$? Давление насыщенного пара $p_{\text{н}2}$ при температуре $t_2$ равно $2{,}5$~кПа. Объёмом жидкости, образовавшейся при конденсации пара, пренебречь по сравнению с $V$.

\kim{24}
\end{taskbasic}
\answer{0{,}29}

% === ЗАДАЧА 2 ===
% Тег: #МКТ #ВЛАЖНОСТЬ #КИМ_24 #ВАРИАНТ_02
\begin{taskbasic}[code={\label{task:mkt_hum_v2_n24}}]
\textbf{Задача 2.} В герметичном сосуде находится водяной пар при температуре $t_1 = 150^{\circ}\text{C}$ и давлении $p_1 = 8$~кПа. Определите объём сосуда, если при охлаждении пара до температуры $t_2 = 20^{\circ}\text{C}$ в нём сконденсируется $\Delta m = 0{,}9$~г воды. Давление насыщенного пара $p_{\text{н}2}$ при температуре $t_2$ равно $2{,}5$~кПа. Объёмом жидкости, образовавшейся при конденсации пара, пренебречь по сравнению с объёмом сосуда.

\kim{24}
\end{taskbasic}
\answer{0{,}9}

% ====================================================================
% ФАЙЛ: tasks/02_mkt/mkt_humidity_bank_part2.tex
% ТЕМА: ВЛАЖНОСТЬ ВОЗДУХА. НАСЫЩЕННЫЕ ПАРЫ
% ПАЧКА 2: ВАРИАНТЫ 7–12
% ====================================================================

% === ЗАДАЧА 1 ===
% Тег: #МКТ #ВЛАЖНОСТЬ #КИМ_22 #ВАРИАНТ_07
\begin{taskbasic}[code={\label{task:mkt_hum_v7_n22}}]
\textbf{Задача 1.} В сосуде содержится влажный воздух при температуре $22^{\circ}\text{C}$. Плотность паров воды в сосуде равна $13\text{ г/м}^3$. Зависимость давления насыщенных паров от температуры представлена на графике. Определите относительную влажность воздуха в сосуде.

\nopagebreak\smallskip
\begin{center}
\begin{tikzpicture}[scale=1.1, every node/.style={font=\footnotesize}]
    \draw[xstep=1, ystep=1, gray!30, thin] (0,0) grid (4,4);
    \draw[-{Stealth[scale=0.8]}, black, thick] (0,0) -- (4.5,0) node[right] {$t, presentation^{\circ}\text{C}$};
    \draw[-{Stealth[scale=0.8]}, black, thick] (0,0) -- (0,4.5) node[above] {$p_{\text{н}}$, гПа};
    
    \draw[ultra thick, brandPrimary] (0,0.4) plot[domain=0:4, samples=20] (\x, {0.4 + 0.15*\x^2 + 0.02*\x^3});
    
    \fill[black] (2.2,1.3) circle (1.5pt) node[above left] {26};
    \draw[dashed, gray] (2.2,0) -- (2.2,1.3) -- (0,1.3);
    
    \node[below] at (2.2,0) {22};
    \node[left] at (0,2) {40};
    \node[left] at (0,3) {60};
    \node[left] at (0,4) {80};
    \node[below] at (1,0) {10};
    \node[below] at (2,0) {20};
    \node[below] at (3,0) {30};
    \node[below] at (4,0) {40};
    \node[below left] at (0,0) {0};
\end{tikzpicture}
\end{center}

\kim{22}
\end{taskbasic}
\answer{50}

% === ЗАДАЧА 2 ===
% Тег: #МКТ #ВЛАЖНОСТЬ #КИМ_22 #ВАРИАНТ_08
\begin{taskbasic}[code={\label{task:mkt_hum_v8_n22}}]
\textbf{Задача 2.} В сосуде содержится влажный воздух при температуре $30^{\circ}\text{C}$. Плотность паров воды в сосуде равна $15\text{ г/м}^3$. Зависимость давления насыщенных паров от температуры представлена на графике. Определите относительную влажность воздуха в сосуде.

\nopagebreak\smallskip
\begin{center}
\begin{tikzpicture}[scale=1.1, every node/.style={font=\footnotesize}]
    \draw[xstep=1, ystep=1, gray!30, thin] (0,0) grid (4,4);
    \draw[-{Stealth[scale=0.8]}, black, thick] (0,0) -- (4.5,0) node[right] {$t, \text{C}$};
    \draw[-{Stealth[scale=0.8]}, black, thick] (0,0) -- (0,4.5) node[above] {$p_{\text{н}}$, гПа};
    
    \draw[ultra thick, brandPrimary] (0,0.4) plot[domain=0:4, samples=20] (\x, {0.4 + 0.15*\x^2 + 0.02*\x^3});
    
    \fill[black] (3.0,2.1) circle (1.5pt) node[above left] {42};
    \draw[dashed, gray] (3.0,0) -- (3.0,2.1) -- (0,2.1);
    
    \node[below] at (3.0,0) {30};
    \node[left] at (0,1) {20};
    \node[left] at (0,2) {40};
    \node[left] at (0,3) {60};
    \node[left] at (0,4) {80};
    \node[below] at (1,0) {10};
    \node[below] at (2,0) {20};
    \node[below] at (3,0) {30};
    \node[below] at (4,0) {40};
    \node[below left] at (0,0) {0};
\end{tikzpicture}
\end{center}

\kim{22}
\end{taskbasic}
\answer{50}

% === ЗАДАЧА 3 ===
% Тег: #МКТ #ВЛАЖНОСТЬ #КИМ_07 #ВАРИАНТ_11
\begin{taskbasic}[code={\label{task:mkt_hum_v11_n7}}]
\textbf{Задача 3.} В закрытом сосуде под поршнем при температуре $t = 100^{\circ}\text{C}$ находится влажный воздух с относительной влажностью $60\%$. Чему станет равна относительная влажность воздуха в сосуде, если при неизменной температуре уменьшить его объём в $1{,}5$ раза?

\kim{7}
\end{taskbasic}
\answer{90}

% === ЗАДАЧА 4 ===
% Тег: #МКТ #ВЛАЖНОСТЬ #КИМ_07 #ВАРИАНТ_12
\begin{taskbasic}[code={\label{task:mkt_hum_v12_n7}}]
\textbf{Задача 4.} В закрытом сосуде под поршнем при постоянной температуре находится влажный воздух с относительной влажностью $40\%$. Чему станет равна относительная влажность воздуха в сосуде, если изотермически уменьшить его объём в $2$ раза?

\kim{7}
\end{taskbasic}
\answer{80}

% ====================================================================
% ФАЙЛ: tasks/02_mkt/mkt_humidity_bank_part3.tex
% ТЕМА: ВЛАЖНОСТЬ ВОЗДУХА. НАСЫЩЕННЫЕ ПАРЫ
% ПАЧКА 3: ВАРИАНТЫ 13–20
% ====================================================================

% === ЗАДАЧА 1 ===
% Тег: #МКТ #ВЛАЖНОСТЬ #КИМ_24 #ВАРИАНТ_15
\begin{taskbasic}[code={\label{task:mkt_hum_v15_n24}}]
\textbf{Задача 1.} В герметичном сосуде объёмом $V = 80$ л находится водяной пар при температуре $t_1 = 150^{\circ}\text{C}$ и давлении $p_1 = 8$ кПа. Какая масса воды $\Delta m$ сконденсируется в сосуде при охлаждении пара до температуры $t_2 = 20^{\circ}\text{C}$? Давление насыщенного пара $p_{\text{н}2}$ при температуре $t_2$ равно 2,5 кПа. Объёмом жидкости, образовавшейся при condensation пара, пренебречь по сравнению с $V$.

\kim{24}
\end{taskbasic}
\answer{0,9}

% === ЗАДАЧА 2 ===
% Тег: #МКТ #ВЛАЖНОСТЬ #КИМ_24 #ВАРИАНТ_16
\begin{taskbasic}[code={\label{task:mkt_hum_v16_n24}}]
\textbf{Задача 2.} В герметичном сосуде находится водяной пар при температуре $t_1 = 150^{\circ}\text{C}$ и давлении $p_1 = 8$ кПа. Определите объём сосуда, если при охлаждении пара до температуры $t_2 = 20^{\circ}\text{C}$ в нём сконденсируется $\Delta m = 0,9$ г воды. Давление насыщенного пара $p_{\text{н}2}$ при температуре $t_2$ равно 2,5 кПа. Объёмом жидкости, образовавшейся при конденсации пара, пренебречь по сравнению с объёмом сосуда.

\kim{24}
\end{taskbasic}
\answer{0,001}

% ====================================================================
% ФАЙЛ: tasks/02_mkt/mkt_humidity_bank_part4.tex
% ТЕМА: ВЛАЖНОСТЬ ВОЗДУХА. НАСЫЩЕННЫЕ ПАРЫ
% ПАЧКА 4: ВАРИАНТЫ 21–30
% ====================================================================

% === ЗАДАЧА 1 ===
% Тег: #МКТ #ВЛАЖНОСТЬ #КИМ_22 #ВАРИАНТ_21
\begin{taskbasic}[code={\label{task:mkt_hum_v21_n22}}]
\textbf{Задача 1.} В сосуде содержится влажный воздух при температуре $22^{\circ}\text{C}$. Плотность паров воды в сосуде равна $13\text{ г/м}^3$. Зависимость давления насыщенных паров от температуры представлена на графике. Определите относительную влажность воздуха в сосуде.

\nopagebreak\smallskip
\begin{center}
\begin{tikzpicture}[scale=1.1, every node/.style={font=\footnotesize}]
    \draw[xstep=1, ystep=1, gray!30, thin] (0,0) grid (4,4);
    \draw[-{Stealth[scale=0.8]}, black, thick] (0,0) -- (4.5,0) node[right] {$t, \,^{\circ}\text{C}$};
    \draw[-{Stealth[scale=0.8]}, black, thick] (0,0) -- (0,4.5) node[above] {$p_{\text{н}}$, гПа};
    
    \draw[ultra thick, brandPrimary] (0,0.4) plot[domain=0:4, samples=20] (\x, {0.4 + 0.15*\x^2 + 0.02*\x^3});
    
    \fill[black] (2.2,1.3) circle (1.5pt) node[above left] {26};
    \draw[dashed, gray] (2.2,0) -- (2.2,1.3) -- (0,1.3);
    
    \node[below] at (2.2,0) {22};
    \node[left] at (0,1) {20};
    \node[left] at (0,2) {40};
    \node[left] at (0,3) {60};
    \node[left] at (0,4) {80};
    \node[below] at (1,0) {10};
    \node[below] at (2,0) {20};
    \node[below] at (3,0) {30};
    \node[below] at (4,0) {40};
    \node[below left] at (0,0) {0};
\end{tikzpicture}
\end{center}

\kim{22}
\end{taskbasic}
\answer{50}

% === ЗАДАЧА 2 ===
% Тег: #МКТ #ВЛАЖНОСТЬ #КИМ_22 #ВАРИАНТ_22
\begin{taskbasic}[code={\label{task:mkt_hum_v22_n22}}]
\textbf{Задача 2.} В сосуде содержится влажный воздух при температуре $30^{\circ}\text{C}$. Плотность паров воды в сосуде равна $15\text{ г/м}^3$. Зависимость давления насыщенных паров от температуры представлена на графике. Определите относительную влажность воздуха в сосуде.

\nopagebreak\smallskip
\begin{center}
\begin{tikzpicture}[scale=1.1, every node/.style={font=\footnotesize}]
    \draw[xstep=1, ystep=1, gray!30, thin] (0,0) grid (4,4);
    \draw[-{Stealth[scale=0.8]}, black, thick] (0,0) -- (4.5,0) node[right] {$t, \,^{\circ}\text{C}$};
    \draw[-{Stealth[scale=0.8]}, black, thick] (0,0) -- (0,4.5) node[above] {$p_{\text{н}}$, гПа};
    
    \draw[ultra thick, brandPrimary] (0,0.4) plot[domain=0:4, samples=20] (\x, {0.4 + 0.15*\x^2 + 0.02*\x^3});
    
    \fill[black] (3.0,2.1) circle (1.5pt) node[above left] {42};
    \draw[dashed, gray] (3.0,0) -- (3.0,2.1) -- (0,2.1);
    
    \node[below] at (3.0,0) {30};
    \node[left] at (0,1) {20};
    \node[left] at (0,2) {40};
    \node[left] at (0,3) {60};
    \node[left] at (0,4) {80};
    \node[below] at (1,0) {10};
    \node[below] at (2,0) {20};
    \node[below] at (3,0) {30};
    \node[below] at (4,0) {40};
    \node[below left] at (0,0) {0};
\end{tikzpicture}
\end{center}

\kim{22}
\end{taskbasic}
\answer{50}